% Part: sets-functions-relations
% Chapter: relations
% Section: orders

\documentclass[../../../include/open-logic-section]{subfiles}

\begin{document}

\olfileid{sfr}{rel}{ord}
\olsection{Tatanan}

\begin{explain}
Akeh pambandhing kita nggambarake objek minangka
"kurang saka", "padha karo", utawa "luwih saka" objek
liyane ing sawijining babagan. Iki nggunakake relasi
\emph{tatanan}. Nanging ana maneka jinis relasi tatanan.
Tuladhane, sawetara mbutuhake saben rong objek bisa
dibandhingake, dene liyane ora.
Sawetara nyakup identitas (kaya~$\le$), lan sawetara
ora nyakup (kaya~$<$).
Panggolongan bakal migunani ing kene.
\end{explain}

\begin{defn}[Pratatanan]
Relasi kang refleksif lan transitif diarani \emph{pratatanan}.
\end{defn}

\begin{defn}[Tatanan parsial]
Pratatanan kang uga antisimetris diarani \emph{tatanan parsial}.
\end{defn}

\begin{defn}[Tatanan linear]\ollabel{def:linearorder}
Tatanan parsial kang uga koneks diarani
\emph{tatanan total} utawa \emph{tatanan linear}.
\end{defn}

\begin{ex}
Saben tatanan linear uga tatanan parsial, lan saben
tatanan parsial uga pratatanan, nanging kosok baline
ora lumaku. Relasi universal ing~$A$ iku pratatanan,
amarga refleksif lan transitif. Nanging yen $A$ duwe
luwih saka siji !!{element}, relasi universal ora
antisimetris, mula dudu tatanan parsial.
\end{ex}

\begin{ex}
Gatekna relasi \emph{ora luwih dawa tinimbang}
$\preccurlyeq$ ing~$\Bin^*$:
$x \preccurlyeq y$ yen lan mung yen
$\len{x} \le \len{y}$.
Iki pratatanan (refleksif lan transitif), malah uga
koneks, nanging dudu tatanan parsial amarga ora
antisimetris. Tuladhane, $01\preccurlyeq 10$ lan
$10 \preccurlyeq 01$, nanging $01 \neq 10$.
\end{ex}

\begin{ex}
Tatanan parsial kang wigati yaiku relasi $\subseteq$
ing himpunan kang ngemot himpunan-himpunan.
Umume iki dudu tatanan linear, amarga yen $a \neq b$
lan kita ngrembug
$\Pow{\{a, b\}} = \{\emptyset, \{a\}, \{b\}, \{a,b\}\}$,
kita weruh yen $\{a\} \nsubseteq \{b\}$,
$\{a\} \neq \{b\}$, lan $\{b\}\nsubseteq \{a\}$.
\end{ex}

\begin{ex}
Relasi \emph{bisa mbagi tanpa sisa} menehi tatanan
parsial kang dudu tatanan linear.
Kanggo wilangan wutuh $n$ lan~$m$,
kita nulis $n \mid m$ kanggo nyatakake yen $n$
mbagi $m$ tanpa sisa, yaiku yen lan mung yen
ana wilangan wutuh~$k$ kang ndadekake $m = kn$.
Ing~$\Nat$, iki tatanan parsial nanging dudu
tatanan linear: tuladhane, $2 \nmid 3$ lan uga
$3 \nmid2$. Yen dianggep relasi ing $\Int$,
relasi bisa mbagi mung pratatanan amarga ora
antisimetris: $1 \mid -1$ lan $-1 \mid 1$,
nanging $1 \neq -1$.
\end{ex}

\begin{ex}
Relasi \emph{ekstensi} ing himpunan urutan~$A^*$
yaiku mangkene: $s \sqsubseteq s'$ yen lan mung yen
$s = \emptyseq$ (urutan kosong), $s = s'$,
utawa $s = \tuple{s_1, \dots, s_n}$ lan
$s' =\tuple{s_1, \dots, s_n, s_{n+1}, \dots, s_m}$.
Yen $s \sqsubseteq s'$, kita uga nyebut $s$
minangka \emph{perangan wiwitan} saka~$s'$.
Relasi ekstensi ing $A^*$ iku tatanan parsial
nanging dudu tatanan linear, tuladhane yen
$a \neq b$, mula $ab \not\sqsubseteq ba$ lan
$ba \not\sqsubseteq ab$.
\end{ex}

\begin{defn}[Tatanan ketat]
\emph{Tatanan ketat} yaiku relasi kang irrefleksif,
asimetris, lan transitif.
\end{defn}

\begin{defn}[Tatanan linear ketat]\ollabel{def:strictlinearorder}
Tatanan ketat kang uga koneks diarani
\emph{tatanan total ketat} utawa \emph{tatanan linear ketat}.
\end{defn}

\begin{ex}
$\le$ iku tatanan linear kang cocog karo
tatanan linear ketat~$<$.
$\subseteq$ iku tatanan parsial kang cocog karo
tatanan ketat~$\subsetneq$.
\end{ex}

Saben tatanan ketat $R$ ing~$A$ bisa diowahi dadi
tatanan parsial kanthi nambah diagonal $\Id{A}$,
yaiku nambah kabeh pasangan~$\tuple{x,x}$.
(Iki diarani \emph{tutupan refleksif} saka~$R$.)
Kosok baline, saka tatanan parsial kita bisa oleh
tatanan ketat kanthi mbuwang~$\Id{A}$.
Rong asil sabanjure njlentrehake iki kanthi persis.

\begin{prop}\ollabel{prop:stricttopartial}
Yen $R$ iku tatanan ketat ing~$A$,
mula $R^+ = R \cup \Id{A}$ iku tatanan parsial.
Kajaba iku, yen $R$ iku tatanan linear ketat,
mula $R^+$ iku tatanan linear.
\end{prop}

\begin{proof}
Umpamakna $R$ iku tatanan ketat, yaiku
$R \subseteq A^2$ lan $R$ irrefleksif,
asimetris, lan transitif.
Tetepna $R^+ = R \cup \Id{A}$.
Kita kudu nuduhake yen $R^+$ refleksif,
antisimetris, lan transitif.

$R^+$ cetha refleksif, amarga
$\tuple{x, x} \in \Id{A} \subseteq R^+$
kanggo kabeh $x \in A$.

Kanggo nuduhake yen $R^+$ antisimetris,
umpamakna kanggo bukti kanthi kontradiksi yen
$R^+xy$ lan $R^+yx$ nanging $x \neq y$.
Amarga $\tuple{x,y} \in R \cup \Id{A}$,
nanging $\tuple{x, y} \notin \Id{A}$,
kita kudu duwe $\tuple{x, y} \in R$, yaiku $Rxy$.
Semono uga,~$Ryx$.
Nanging iki nalisir anggepan yen $R$ asimetris.

Kanggo mbuktekake transitivitas, umpamakna
$R^+xy$ lan $R^+yz$.
Yen $\tuple{x, y} \in R$ lan $\tuple{y,z} \in R$
loro-lorone lumaku, mula $\tuple{x, z} \in R$
amarga $R$~transitif.
Yen ora, mula $\tuple{x, y} \in\Id{A}$,
yaiku $x = y$, utawa $\tuple{y, z} \in \Id{A}$,
yaiku $y = z$.
Ing kasus kapisan, kita duwe $R^+yz$ miturut
anggepan lan $x = y$, mula $R^+xz$.
Semono uga ing kasus kapindho.
Ing saben kasus, $R^+xz$, mula $R^+$
uga transitif.

Tumrap pratelan "kajaba iku", umpamakna $R$
uga koneks. Mula kanggo kabeh $x \neq y$,
$Rxy$ utawa~$Ryx$, yaiku
$\tuple{x, y} \in R$ utawa $\tuple{y, x} \in R$.
Amarga $R \subseteq R^+$, iki tetep lumaku
tumrap $R^+$, mula $R^+$ uga koneks.
\end{proof}

\begin{prop}\ollabel{prop:partialtostrict}
Yen $R$ iku tatanan parsial ing~$A$,
mula $R^- = R \setminus \Id{A}$ iku tatanan ketat.
Kajaba iku, yen $R$ iku tatanan linear,
mula $R^-$ iku tatanan linear ketat.
\end{prop}

\begin{proof}
Iki ditinggalake minangka latihan.
\end{proof}

\begin{prob}
Wenehana bukti kanggo \olref[sfr][rel][ord]{prop:partialtostrict}.
\end{prob}

Asil prasaja ing ngisor iki nuduhake yen
tatanan linear ketat nduweni sipat kang
padha karo ekstensionalitas:

\begin{prop}\ollabel{prop:extensionality-strictlinearorders}
Yen $<$ iku tatanan linear ketat ing $A$, mula:
\[
  (\forall a, b \in A)((\forall x \in A)(x < a \liff x < b) \lif a = b).
\]
\end{prop}

\begin{proof}
Umpamakna $(\forall x \in A)(x < a \liff x < b)$.
Yen $a < b$, mula $a <a$,
kang nalisir kasunyatan yen $<$ irrefleksif;
mula $a \nless b$.
Kanthi cara kang persis padha, $b \nless a$.
Mula $a = b$, amarga $<$ koneks.
\end{proof}

\end{document}
